\documentclass{article}%
\usepackage[T1]{fontenc}%
\usepackage[utf8]{inputenc}%
\usepackage{lmodern}%
\usepackage{textcomp}%
\usepackage{lastpage}%
\usepackage{geometry}%
\geometry{margin=1in,headheight=14pt,headsep=25pt}%
\usepackage{hyperref}%
\usepackage{enumitem}%
\usepackage{fancyhdr}%
\usepackage{xcolor}%
\usepackage{url}%
\usepackage{breakurl}%
%

            \hypersetup{
                colorlinks=true,
                linkcolor=blue,
                filecolor=magenta,
                urlcolor=blue
            }
            \pagestyle{fancy}
            \fancyhf{}
            \rhead{Generated on \today}
            \cfoot{\thepage}
            
            % Configure enumeration settings
            \setlist[enumerate,1]{label=\arabic*.}
            \setlist[enumerate,2]{label=\alph*.}
            \setlist[enumerate,3]{label=\Alph*.}
            \setlist[enumerate]{itemsep=0.5em}
            
            % Define a command for red text
            \newcommand{\incorrect}[1]{\textcolor{red}{#1}}
        %
\lhead{\$L\^{}\textbackslash{}infty\${-}Regression}%
\title{Practice Questions for \$L\^{}\textbackslash{}infty\${-}Regression}%
\author{Generated by Scribe.AI}%
\date{\today}%
%
\begin{document}%
\normalsize%
\maketitle%
\section{Questions}%
\label{sec:Questions}%
\begin{enumerate}%
\item In $L^\infty$-regression, what is the primary objective of the optimization problem, and how does this objective differ from that of $L^2$-regression and $L^1$-regression?%
\begin{enumerate}%
\item To minimize the sum of absolute errors between the predicted and actual values, similar to $L^1$-regression.%
\item To minimize the maximum absolute error between the predicted and actual values.%
\item To minimize the sum of squared errors between the predicted and actual values, similar to $L^2$-regression.%
\item To minimize the average absolute error between the predicted and actual values.%
\item To maximize the minimum absolute error between the predicted and actual values.%
\end{enumerate}%
\vspace{1em}%
\item How does the choice of the loss function ($L^\infty$ vs. $L^2$ vs. $L^1$) affect the sensitivity of the regression model to outliers in the dataset?  Discuss the relative robustness of each method.%
\begin{enumerate}%
\item $L^\infty$ is most sensitive to outliers, followed by $L^2$, and $L^1$ is least sensitive.%
\item $L^1$ is most sensitive to outliers, followed by $L^\infty$, and $L^2$ is least sensitive.%
\item $L^2$ is most sensitive to outliers, followed by $L^1$, and $L^\infty$ is least sensitive.%
\item $L^\infty$ is least sensitive to outliers, followed by $L^1$, and $L^2$ is most sensitive.%
\item All three methods are equally sensitive to outliers.%
\end{enumerate}%
\vspace{1em}%
\item What is a key computational challenge associated with solving the $L^\infty$-regression problem compared to $L^2$-regression, and why does this challenge arise?%
\begin{enumerate}%
\item $L^\infty$-regression is computationally simpler than $L^2$-regression because it involves only linear equations.%
\item $L^\infty$-regression often requires linear programming techniques, which can be more computationally intensive than the closed-form solution available for $L^2$-regression.%
\item There is no significant computational difference between solving $L^\infty$ and $L^2$ regression problems.%
\item $L^\infty$-regression is always faster because it only considers the maximum error.%
\item $L^2$-regression is computationally more challenging due to the need for iterative methods.%
\end{enumerate}%
\vspace{1em}%
\item Describe the geometric interpretation of the solution to the $L^\infty$-regression problem, and how does it differ from the geometric interpretation of the $L^2$-regression solution?%
\begin{enumerate}%
\item Both $L^2$ and $L^\infty$ regression solutions minimize the sum of squared distances between data points and the fitted line.%
\item $L^2$-regression finds the line that minimizes the sum of squared distances, while $L^\infty$-regression finds the line that minimizes the maximum vertical distance between data points and the fitted line.%
\item $L^2$-regression finds the line that minimizes the maximum vertical distance, while $L^\infty$-regression minimizes the sum of squared distances.%
\item There is no geometric interpretation for $L^\infty$-regression.%
\item Both methods find the line that minimizes the sum of absolute distances between data points and the fitted line.%
\end{enumerate}%
\vspace{1em}%
\item In $L^\infty$-regression, what is the primary objective of the optimization problem?%
\begin{enumerate}%
\item Minimize the sum of absolute errors.%
\item Minimize the maximum absolute error.%
\item Minimize the sum of squared errors.%
\item Minimize the average absolute error.%
\item Minimize the median absolute error.%
\end{enumerate}%
\vspace{1em}%
\item How does the choice of the loss function ($L^\infty$ vs. $L^2$ vs. $L^1$) affect the sensitivity of the regression model to outliers?%
\begin{enumerate}%
\item $L^\infty$ is most sensitive, $L^2$ is intermediate, and $L^1$ is least sensitive to outliers.%
\item $L^1$ is most sensitive, $L^2$ is intermediate, and $L^\infty$ is least sensitive to outliers.%
\item $L^2$ is most sensitive, $L^1$ is intermediate, and $L^\infty$ is least sensitive to outliers.%
\item $L^\infty$ is least sensitive, $L^2$ is most sensitive, and $L^1$ is intermediate.%
\item $L^1$ is least sensitive, $L^2$ is most sensitive, and $L^\infty$ is intermediate.%
\end{enumerate}%
\vspace{1em}%
\item What is a key computational challenge associated with solving the $L^\infty$-regression problem compared to $L^2$-regression?%
\begin{enumerate}%
\item The $L^\infty$ problem is always easier to solve than the $L^2$ problem.%
\item The $L^\infty$ problem is a linear program, while the $L^2$ problem is not.%
\item The $L^\infty$ problem often requires iterative methods, unlike the $L^2$ problem which has a closed-form solution.%
\item The $L^\infty$ problem is always computationally intractable, while the $L^2$ problem is always tractable.%
\item There is no significant computational difference between solving $L^\infty$ and $L^2$ regression problems.%
\end{enumerate}%
\vspace{1em}%
\item Describe the geometric interpretation of the solution to the $L^\infty$-regression problem, and how does it differ from the geometric interpretation of the $L^2$-regression solution?%
\begin{enumerate}%
\item Both minimize the sum of squared distances to the fitted line.%
\item $L^2$ minimizes the sum of squared distances, while $L^\infty$ minimizes the maximum distance.%
\item $L^2$ minimizes the maximum distance, while $L^\infty$ minimizes the sum of squared distances.%
\item Both minimize the maximum distance to the fitted line.%
\item There is no geometric interpretation for $L^\infty$ regression.%
\end{enumerate}%
\vspace{1em}%
\end{enumerate}

%
\newpage%
\section{Answers}%
\label{sec:Answers}%
\begin{enumerate}%
\item In $L^\infty$-regression, what is the primary objective of the optimization problem, and how does this objective differ from that of $L^2$-regression and $L^1$-regression?%
\begin{enumerate}%
\item \incorrect{This describes the objective of $L^1$-regression, not $L^\infty$-regression.}%
\item The primary objective in $L^\infty$-regression is to minimize the maximum absolute error. This differs from $L^2$-regression, which minimizes the sum of squared errors, and $L^1$-regression, which minimizes the sum of absolute errors.%
\item \incorrect{This describes the objective of $L^2$-regression, not $L^\infty$-regression.}%
\item \incorrect{Minimizing the average absolute error is a different objective, not directly related to $L^\infty$-regression.}%
\item \incorrect{Maximizing the minimum error is not the goal; the aim is to minimize the maximum error.}%
\end{enumerate}%
\vspace{1em}%
\item How does the choice of the loss function ($L^\infty$ vs. $L^2$ vs. $L^1$) affect the sensitivity of the regression model to outliers in the dataset?  Discuss the relative robustness of each method.%
\begin{enumerate}%
\item \incorrect{This is incorrect; $L^1$ is the most robust to outliers.}%
\item \incorrect{This is incorrect; $L^2$ is more sensitive to outliers than $L^1$ and $L^\infty$.}%
\item \incorrect{This is incorrect; $L^\infty$ is the least sensitive to outliers.}%
\item The $L^\infty$ norm focuses on the single largest error, making it relatively insensitive to outliers compared to $L^1$ and $L^2$. $L^1$ is more robust to outliers than $L^2$ because it only considers the magnitude of errors, not their squares.%
\item \incorrect{The different norms have different sensitivities to outliers.}%
\end{enumerate}%
\vspace{1em}%
\item What is a key computational challenge associated with solving the $L^\infty$-regression problem compared to $L^2$-regression, and why does this challenge arise?%
\begin{enumerate}%
\item \incorrect{This is incorrect; $L^\infty$-regression is generally more computationally intensive.}%
\item Unlike $L^2$-regression, which has a closed-form solution via normal equations, $L^\infty$-regression often requires linear programming or other iterative methods, making it computationally more demanding.%
\item \incorrect{There is a significant computational difference; $L^2$ has a closed-form solution, while $L^\infty$ often requires iterative methods.}%
\item \incorrect{This is incorrect; the iterative nature of solving $L^\infty$ problems makes them slower.}%
\item \incorrect{This is incorrect; $L^2$ has a closed-form solution.}%
\end{enumerate}%
\vspace{1em}%
\item Describe the geometric interpretation of the solution to the $L^\infty$-regression problem, and how does it differ from the geometric interpretation of the $L^2$-regression solution?%
\begin{enumerate}%
\item \incorrect{This is incorrect; $L^\infty$ minimizes the maximum distance, not the sum of squared distances.}%
\item The $L^2$ solution minimizes the sum of squared vertical distances, while the $L^\infty$ solution minimizes the maximum vertical distance between data points and the fitted line.  This leads to different sensitivities to outliers.%
\item \incorrect{This is incorrect; it reverses the geometric interpretations of the two methods.}%
\item \incorrect{This is incorrect; $L^\infty$-regression has a geometric interpretation.}%
\item \incorrect{This is incorrect; this describes $L^1$-regression.}%
\end{enumerate}%
\vspace{1em}%
\item In $L^\infty$-regression, what is the primary objective of the optimization problem?%
\begin{enumerate}%
\item \incorrect{This describes the objective of $L^1$-regression.}%
\item The primary objective in $L^\infty$-regression is to minimize the maximum absolute error between the predicted and actual values. This is in contrast to $L^1$ and $L^2$ regression, which minimize the sum of absolute errors and the sum of squared errors, respectively.%
\item \incorrect{This describes the objective of $L^2$-regression.}%
\item \incorrect{Minimizing the average absolute error is related to $L^1$-regression but is not the primary focus of $L^\infty$-regression.}%
\item \incorrect{Minimizing the median absolute error is a robust regression technique but is not the same as $L^\infty$-regression.}%
\end{enumerate}%
\vspace{1em}%
\item How does the choice of the loss function ($L^\infty$ vs. $L^2$ vs. $L^1$) affect the sensitivity of the regression model to outliers?%
\begin{enumerate}%
\item \incorrect{Incorrect ordering of sensitivity to outliers.}%
\item \incorrect{Incorrect ordering of sensitivity to outliers.}%
\item $L^2$ regression is most sensitive to outliers because squaring large errors amplifies their influence. $L^1$ regression is less sensitive because it uses absolute values, and $L^\infty$ regression is the least sensitive because it only considers the maximum error, making it robust to outliers.%
\item \incorrect{Incorrect ordering of sensitivity to outliers.}%
\item \incorrect{Incorrect ordering of sensitivity to outliers.}%
\end{enumerate}%
\vspace{1em}%
\item What is a key computational challenge associated with solving the $L^\infty$-regression problem compared to $L^2$-regression?%
\begin{enumerate}%
\item \incorrect{Incorrect; $L^\infty$ regression is generally more computationally intensive.}%
\item \incorrect{Both can be formulated as linear programs, but the $L^2$ problem has a simpler, closed-form solution.}%
\item The $L^2$ problem has a closed-form solution via the normal equations.  The $L^\infty$ problem is typically solved using iterative methods, which can be computationally more expensive, especially for large datasets.%
\item \incorrect{Incorrect; both problems are generally tractable, but $L^\infty$ is often more computationally expensive.}%
\item \incorrect{Incorrect; there is a significant computational difference.}%
\end{enumerate}%
\vspace{1em}%
\item Describe the geometric interpretation of the solution to the $L^\infty$-regression problem, and how does it differ from the geometric interpretation of the $L^2$-regression solution?%
\begin{enumerate}%
\item \incorrect{Incorrect; $L^2$ minimizes the sum of squared distances, but $L^\infty$ minimizes the maximum distance.}%
\item In $L^2$ regression, the solution minimizes the sum of squared vertical distances from the data points to the fitted line. In $L^\infty$ regression, the solution minimizes the maximum vertical distance from any data point to the fitted line.  This means the fitted line is contained within a band of width equal to twice the minimum maximum distance.%
\item \incorrect{Incorrect; the descriptions of $L^2$ and $L^\infty$ are reversed.}%
\item \incorrect{Incorrect; only $L^\infty$ minimizes the maximum distance.}%
\item \incorrect{Incorrect; $L^\infty$ regression has a geometric interpretation.}%
\end{enumerate}%
\vspace{1em}%
\end{enumerate}

%
\end{document}